\documentclass[11pt, a4paper]{common/document_template}
\usepackage{common}

\begin{document}

\begin{titlepage}
  \begin{center}
	\Huge{Турнир Савина}
  \end{center}
\end{titlepage} 

\chapter{2025. Номер 30}

\begin{exercise}
По кругу стоят 100 детей. Каждый из них подсчитал, сколько детей
противоположного пола среди следующих по часовой стрелке десяти детей. Докажите,
что суммы чисел, полученных мальчиками и девочками, равны.
\end{exercise}

\begin{exercise}[13]
В каждой клетке доски $9 \times 9$ сидел муравей. В некоторый момент каждый муравей
переполз в соседнюю по стороне клетку. Могло ли после этого в каждой непустой клетке
оказаться ровно по три муравья?
\end{exercise}

\begin{Solution}
Раскрасим доску шахматной раскраской. Изначально суммарное количество муравьев в черных клетках - 41, а в белых - 40.
То есть, разница между черными и белыми клетками равна 1. Так как при переползании каждый муравей меняет цвет клетки, 
то эта разница по модулю не меняется. Но, с другой стороны, эта разница по условию должна делиться на 3. Значит, по три муравья
в каждой непустой клетке оказаться не могло.
\end{Solution}

\begin{exercise}[31]
В компании 50 человек, некоторые из них дружат друг с другом. Людей
пронумеровали различными числами от 1 до 50, при этом оказалось, что номер каждого
человека является делителем суммы номеров его друзей. Докажите, что у кого-то в этой
компании чётное число друзей.
\end{exercise}

\begin{exercise}
Назовём натуральное число почётным, если в его записи сумма нечётных цифр
равна сумме чётных. Докажите, что для любого натурального $N$ найдётся почётное число,
кратное $N$.
\end{exercise}

\begin{exercise}[36]
На плоскости провели 20 прямых. Какое наибольшее число прямоугольников может
быть образовано этими прямыми?
\end{exercise}

\begin{exercise}[86]
Из квадрата, разбитого на единичные клетки, вырезали клетчатый многоугольник
площади 21 и периметра 42. Какой наименьшей длины могла быть сторона квадрата?
\end{exercise}

\begin{exercise}[106]
В комнате находятся десять человек. У одного из них ровно четыре знакомых в
этой комнате, а у остальных - больше. Докажите, что найдётся тройка человек, в которой
каждый знаком с двумя остальными.
\end{exercise}

\begin{Solution}
Отнесем этих четырех знакомых в группу $A$. А оставшихся 5 человек в группу $B$. 
Тогда если в группе $A$ есть знакомые между собой, то утверждение доказано.
Предположим, что в группе $A$ знакомых между собой нет. Тогда у каждого из группы $A$ не менее четырех знакомых в группе $B$.
В группе $B$ пять человек на которых приходится 16 знакомств из группы $A$. Значит, по принципу Дирихле, 
в группе $B$ есть человек, который знаком со всеми из группы $A$. Но, исходя из условия, он должен быть знаком с кем то из группы $B$.
А значит образуется тройка из двух человек в группе $B$ и одного из $A$, в  которой каждый знаком с двумя остальными.
\end{Solution}

\begin{exercise}[118]
Существуют ли два натуральных числа, отличающихся ровно в 7 раз, десятичные
записи которых получаются друг из друга циклическим сдвигом?
\end{exercise}

\begin{exercise}[123]
Назовём натуральное число $N$ особенным, если в его десятичной записи
встречаются только две различные ненулевые цифры, а число $9N$ представимо в виде
разности кубов двух натуральных чисел. Докажите, что особенных чисел бесконечно
много.
\end{exercise}

\begin{exercise}[125]
Можно ли разрезать треугольник на выпуклые пятиугольники?
\end{exercise}

\begin{exercise}[129]
Докажите, что уравнение
$$
\frac{1}{x^{2}} + \frac{1}{y^{2}} + \frac{1}{z^{2}} = \frac{100}{x^{2} + y^{2} + z^{2}}
$$
не имеет решений в целых числах.
\end{exercise}

\end{document}